\subsection{Presentación del caso}

Aquí viene el enunciado del pptx que compartiré contigo


\newpage
\subsection{Formulación explíctia del modelo}

\paragraph{Variables}\mbox{}\\

\begin{tabular}{l c l}
    $x_b$  & : & número de unidades producidas de baguettes\\
    $x_p$  & : & número de unidades producidas de pan de picos\\
    $x_g$  & : & número de unidades producidas de pagn gallego\\
\end{tabular}

\paragraph{Restricciones}\mbox{}\\
\\
No se producen más unidades de las que se pueden vender:

\begin{equation}
    x_b + x_p + x_g \leq 400
\end{equation}
\\
\\
No se pueden producir más productos que los que permite la cantidad máxima de harina disponbible

\begin{equation}
    0.1x_b + 0.2x_p + 0.4x_g \leq 600
\end{equation}
\\
\\
El tiempo necesario para amasar la producción no puede superar el tiempo disponible para amasar:


\begin{equation}
    2x_b + 5x_p + 7x_g \leq 4\times4\times 60 = 960
\end{equation}
\\
\\
\paragraph{Función objetivo}\mbox{}\\
\\
\\
Cómputo de los ingresos:
\begin{equation}
1x_b + 1.5x_p + 2.5x_g
\end{equation}
\\
\\
Cómputo de los gastos en harina
\begin{equation}
(0.1x_b + 0.2x_p + 0.4x_g)\times0.5
\end{equation}
\\
\\
Costes fijos: los salarios de los emplados. Son costes irrelevantes para decidir la producción.
\\
\\
Expresión de la función objetivo:
\begin{equation}
\begin{split}
1x_b + 1.5x_p + 2.5x_g - (0.1x_b + 0.2x_p + 0.4x_g)\times0.5 \\
= 0.95x_b + 1.4x_p + 2.3x_g 
\end{split}
\end{equation}

\paragraph{Modelo completo}\mbox{}\\
\\
\\
El modelo completo es el siguiente:
\begin{align}
\mbox{max.} \quad & 0.95x_b + 1.4x_p + 2.3x_g \notag \\
\mbox{s. a.:} \quad &  \notag \\		
             & x_b + x_p + x_g \leq 400 \notag \\
             &  0.1x_b + 0.2x_p + 0.4x_g \leq 600\notag \\
             & 2x_b + 5x_p + 7x_g \leq 960\notag  \\
             & x_b, x_p. x_g \geq 0
\end{align}    

\subsection{Formulación generalizada}

\paragraph{Conjuntos}\mbox{}\\

\begin{tabular}{l c l}
    $\mathcal{T}$  & : & tipos de panes \\
\end{tabular}
\\
\\

\paragraph{Parámetros}\mbox{}\\

\begin{tabular}{l c l}
    $P$  & : & parámetro 1 \\
    $Q$  & : & parámetro 2 \\
\end{tabular}
\\
\\

\paragraph{Variables}\mbox{}\\

\begin{tabular}{l c l}
    $x_t$  & : & variable 1 para $t \in \mathcal{T}$\\
    $y_t$  & : & variable 1 para $t \in \mathcal{T}$\\
\end{tabular}

\paragraph{Restricciones}\mbox{}\\

Descripción del primer requisito:

\begin{equation}
    \sum_t{x_t} \leq V
\end{equation}

\
\\
Descripción del segundo requisito:

\begin{equation}
    \sum_t{NEC_tx_t} \leq H
\end{equation}
En este caso: 
\\

\paragraph{Función objetivo}\mbox{}\\

\begin{align}
             & \sum_tP_t{x_t}  \qquad \mbox{(ingresos por las ventas)} \notag \\
             & -C\sum_t  NEC_t {x_t} \qquad \mbox{(costes de la harina)} \notag \\
             & -\mbox{costes fijos} \qquad \mbox{(costes de los operarios)} \notag \\
\end{align} 



\paragraph{Modelo completo}\mbox{}\\
El modelo de forma generalizada es el siguiente.
\\
\\
\begin{align}
\mbox{max.} \quad & \sum_tP_t{x_t}  \qquad \mbox{(ingresos por las ventas)} \notag \\
             & -C\sum_t  NEC_t {x_t} \qquad \mbox{(costes de la harina)} \notag \\
             & -\mbox{costes fijos} \qquad \mbox{(costes de los operarios)} \notag \\
\mbox{s. a.:} \quad &  \notag \\		
             & \sum_t{x_t} \leq V \qquad\mbox{(restricciones de mercado)} \notag \\
             &  \sum_t{NEC_tx_t} \leq H \qquad\mbox{(restricciones de materia prima)}   \notag \\
             & \sum_t{AM_tx_t} \leq T \qquad \mbox{(restricciones amasado)} \notag \\
             & x_t,  \geq 0 \; \forall t \in \mathcal{T} \qquad \mbox{(restricciones no negatividad)} \notag
\end{align} 
